\section{Metrics for Loop Engineering Evaluation}\label{app:metrics}

We formalize the five headline RQ1 metrics used to evaluate model and loop configurations on \lhb{}. Together, they separate end state scoring from loop trace metrics for dependency progress, obligation retention, and resource use. Let a task have tests indexed by $i$ and evaluation checkpoints indexed by $\tau \in \{1,\dots,T\}$. Let $s_i^{(\tau)} \in \{0,1\}$ denote the pass/fail outcome of test $i$ at checkpoint $\tau$, and let $O_\tau$ denote the set of \emph{regression obligations} active at checkpoint $\tau$ (tests previously released and passed at some earlier checkpoint).
The recorded test state matrix is $S = [s_i^{(\tau)}]$ and serves as the evidence basis from which both final outcomes and intermediate loop trace metrics are computed.
Let $D_\tau$ denote the set of tests released by checkpoint $\tau$, and let $T_{\max}$ denote the wall clock time budget. The released test pass rate at checkpoint $\tau$ is
\begin{equation}
\operatorname{TPR}(\tau) \;=\; \frac{\bigl|\{\, i \in D_\tau : s_i^{(\tau)} = 1 \,\}\bigr|}{|D_\tau|}.
\end{equation}

\subsection{Resolve Rate (RR)}\label{app:final-success}
\begin{equation}
\operatorname{RR} \;=\; \mathbb{1}\!\left[\,\bigwedge_{i \in \text{all tests}} s_i^{(T)} = 1\,\right].
\end{equation}
RR records whether every released test passes at task end, including regression obligations accumulated as tests are released. It is the strict end state outcome metric, while the remaining metrics record dependency progress, obligation retention, and resource use along the recorded loop trace.

\subsection{Test Pass Rate (TPR)}\label{app:pass-rate}
\begin{equation}
\operatorname{TPR} \;=\; \frac{1}{|D_T|}\sum_{i \in D_T} s_i^{(T)}.
\end{equation}
TPR records the average fraction of released tests passing at task end. Unlike RR, it preserves partial progress over the released test set and connects the final state to the intermediate loop trace recorded during execution.

\subsection{Dependency Depth (Depth)}\label{app:depth}

Depth measures dependency reach before the first residual handoff by the outer evaluation loop, isolating how far one continuous execution segment preserves prerequisite obligations before residual continuation begins. Let each released test $i$ inherit the topological layer of the development unit it adjudicates, $\mathrm{layer}(i) \in \{0, \dots, L\}$, where $L = \max_{v \in V} \mathrm{layer}(v)$ is the longest prerequisite chain of the task DAG. Let $\tau^{\dagger}$ denote the last checkpoint of the first outer evaluation round, that is, the final snapshot before the outer loop, if any, restarts the loop with a freshly derived residual. Let $P_{\tau^{\dagger}} = \{i : s_i^{(\tau^{\dagger})} = 1\}$ denote the tests passing at $\tau^{\dagger}$. Depth is the deepest layer at which every prerequisite test still passes at $\tau^{\dagger}$, normalized by $L$:
\begin{equation}
\resizebox{.9\columnwidth}{!}{$\displaystyle
\operatorname{Depth} \;=\; \mathbb{E}_{\text{task}}\!\left[\,\frac{\max\bigl\{\ell : \forall i,\, \mathrm{layer}(i) \leq \ell \Rightarrow i \in P_{\tau^{\dagger}}\bigr\}}{L}\,\right].
$}
\end{equation}
Depth records dependency reach before the first residual handoff in the evaluated loop, while RR combines this reach with residual recovery across later rounds. Resolved tasks may reach $\mathrm{RR}=1$ through high Depth before the first residual handoff or through lower initial Depth followed by successful residual recovery. Read together, the two metrics separate entry layer stalls ($\operatorname{Depth}{\ll}1$) from late prerequisite discontinuities, where high Depth pairs with low RR after later execution no longer preserves prior obligations.

\subsection{Regression Rate (Reg)}\label{app:regression-rate}

A \emph{regression event} occurs whenever a previously passing obligation test fails at a later evaluation point, making obligation retention observable as a loop trace metric rather than solely as final nonresolution. Checkpoints $\tau$ are snapshots emitted by the dual container watcher (Section~\ref{sec:protocol}). When the cumulative diff from the evaluated loop exceeds a fixed line threshold against the snapshot baseline, the watcher container freezes the workspace and runs the full released test suite, recording per test pass or fail outcomes in \texttt{observation\_history.jsonl} alongside the \texttt{newly\_passed} and \texttt{newly\_failed} deltas relative to the previous snapshot. Let $\tau^-$ denote the immediately preceding triggered snapshot at which test $i$ was evaluated, let $O_\tau$ denote the obligation set at $\tau$ (the union of all tests that were \texttt{newly\_passed} at any earlier triggered snapshot), and define
\begin{equation}
R = \#\bigl\{(i,\tau): i \in O_\tau,\; s_i^{(\tau^-)} \!=\! 1,\; s_i^{(\tau)} \!=\! 0\bigr\}.
\end{equation}
Then
\begin{equation}
\operatorname{Reg} \;=\; \frac{R}{\sum_\tau |O_\tau|}.
\end{equation}
Equivalently, Reg measures how often a previously satisfied obligation stops passing within the recorded loop trace, making obligation retention a directly measured loop trace metric.
Reg is interpreted jointly with Resolve Rate. Configurations with very low Resolve Rate satisfy a limited obligation set, and $|O_\tau|$ therefore provides limited exposure for estimating retention. Low Reg in that regime provides limited evidence of stable long horizon execution. Conversely, a model and loop configuration with moderate RR but high Reg, such as DeepSeek-V4P at $\text{RR}{=}11.61\%$ and $\text{Reg}{=}16.60\%$ in our runs, corresponds to obligation retention degradation beyond limited progress alone. We report Reg as a loop trace metric and use Resolve Rate for outcome ranking.

\subsection{Tokens (Tok)}\label{app:tokens}

Tok records total tokens consumed per task, summed across input and output over the full loop trace. This metric reports the resource cost of sustaining the evaluated loop over long horizon execution and contextualizes loop trace improvements against runtime expenditure.
